\documentclass[reqno]{amsart}
\usepackage[T1]{fontenc}
\usepackage{amsmath,amssymb,amsthm}
\usepackage{xr-hyper}
\usepackage[hidelinks]{hyperref}
\externaldocument[ec-]{20261124T223000--elliptic-curves-for-the-proof__workbook}
\externaldocument[da-]{20261126T030300--the-data__workbook}
\externaldocument[ob-]{20261126T040100--spaces-forms-and-twists__workbook}
\externaldocument[dp-]{20261126T030900--denominators-of-projections-to-lifts__workbook}
\externaldocument[tw-]{20261126T031000--twists-with-few-congruences__workbook}
\externaldocument[c3-]{20261126T050600--the-exchange-identity__workbook}
\externaldocument[l3-]{20261126T050900--hecke-values-as-norms-of-lifts__workbook}
\externaldocument[w2-]{20261126T060600--lengths-as-values-of-polynomials__workbook}
\externaldocument[w3-]{20261126T060700--tied-primes__workbook}

\newcounter{chap}
% Chapter numbers, generated from 20261125T012000--grh-workbook__guide.org
% by scripts/make-chapters.py.  Do not edit: change the order of the
% index and rerun the script (./preview-site.sh and CI do this).
%
% Usage in a sheet, after \newcounter{chap}:
%   % Chapter numbers, generated from 20261125T012000--grh-workbook__guide.org
% by scripts/make-chapters.py.  Do not edit: change the order of the
% index and rerun the script (./preview-site.sh and CI do this).
%
% Usage in a sheet, after \newcounter{chap}:
%   % Chapter numbers, generated from 20261125T012000--grh-workbook__guide.org
% by scripts/make-chapters.py.  Do not edit: change the order of the
% index and rerun the script (./preview-site.sh and CI do this).
%
% Usage in a sheet, after \newcounter{chap}:
%   \input{chapters}\chapnumber{<sheet ID>}
\makeatletter
\providecommand\chapters@set[2]{\@namedef{chapters@#1}{#2}}
\chapters@set{20261125T001000}{1}% zeros-to-ramp-sums
\chapters@set{20261125T001100}{2}% inert-primes-and-estimate-b
\chapters@set{20261125T005000}{3}% galois-fourier-and-averaging
\chapters@set{20261125T005500}{4}% columns-from-a-sphere
\chapters@set{20261125T000000}{5}% slope-method
\chapters@set{20261126T030100}{6}% a-twist-for-each-prime
\chapters@set{20261125T003000}{7}% the-chain-and-the-square-root
\chapters@set{20261125T010000}{8}% rank-lemma-and-zero-estimate
\chapters@set{20261125T010500}{9}% poincare-series-and-near-orthogonality
\chapters@set{20261124T223000}{10}% elliptic-curves-for-the-proof
\chapters@set{20261125T009500}{11}% fock-model-and-theta-intertwiners
\chapters@set{20261124T223100}{12}% complex-multiplication-and-inert-reduction
\chapters@set{20261125T008000}{13}% one-cm-determinant
\chapters@set{20261124T214500}{14}% supersingular-theta-divisibility
\chapters@set{20261124T214600}{15}% the-isotropic-plane
\chapters@set{20261125T013700}{16}% pairings-on-a-formal-group-of-height-two
\chapters@set{20261125T013500}{17}% hecke-nondegeneracy-on-the-tree
\chapters@set{20261126T030200}{18}% the-bound-is-attained
\chapters@set{20261125T009000}{19}% kronecker-limit-and-chowla-selberg
\chapters@set{20261125T002000}{20}% definite-theta-lifts
\chapters@set{20261125T015500}{21}% siegel-weil-in-the-smallest-case
\chapters@set{20261125T002100}{22}% definite-theta-lifts-and-norms
\chapters@set{20261125T016500}{23}% hecke-operators-and-theta-lifts
\chapters@set{20261125T016700}{24}% divisibility-along-the-chain
\chapters@set{20261125T004000}{25}% denominators-and-sha
\chapters@set{20261125T004100}{26}% tate-twists-and-sha
\chapters@set{20261126T020430}{27}% the-minus-part-by-gauss-sums
\chapters@set{20261125T014100}{28}% congruence-modules
\chapters@set{20261125T014500}{29}% eisenstein-congruences-and-selmer-classes
\chapters@set{20261125T014700}{30}% pseudocharacters-and-the-length-bound
\chapters@set{20261125T011000}{31}% galois-representations-black-box
\chapters@set{20261126T030300}{32}% the-data
\chapters@set{20261126T020210}{33}% bessel-poincare-series-and-their-gram-operators
\chapters@set{20261126T030400}{34}% the-seesaw-identity
\chapters@set{20261126T030500}{35}% divisibility-of-the-composite-lift
\chapters@set{20261126T030600}{36}% theorem-a-short-forms
\chapters@set{20261126T030700}{37}% the-determinant-of-a-chain
\chapters@set{20261126T030800}{38}% the-cost-of-one-lift
\chapters@set{20261126T030900}{39}% denominators-of-projections-to-lifts
\chapters@set{20261126T031000}{40}% twists-with-few-congruences
\chapters@set{20261126T031100}{41}% theorem-c-and-the-main-inequality
\providecommand\chapnumber[1]{%
  \@ifundefined{chapters@#1}%
    {\PackageError{chapters}{Sheet #1 is not listed in the index}%
      {List it in 20261125T012000--grh-workbook__guide.org\MessageBreak
       and run scripts/make-chapters.py.}}%
    {\edef\thechap{\@nameuse{chapters@#1}}}}
\makeatother
\chapnumber{<sheet ID>}
\makeatletter
\providecommand\chapters@set[2]{\@namedef{chapters@#1}{#2}}
\chapters@set{20261125T001000}{1}% zeros-to-ramp-sums
\chapters@set{20261125T001100}{2}% inert-primes-and-estimate-b
\chapters@set{20261125T005000}{3}% galois-fourier-and-averaging
\chapters@set{20261125T005500}{4}% columns-from-a-sphere
\chapters@set{20261125T000000}{5}% slope-method
\chapters@set{20261126T030100}{6}% a-twist-for-each-prime
\chapters@set{20261125T003000}{7}% the-chain-and-the-square-root
\chapters@set{20261125T010000}{8}% rank-lemma-and-zero-estimate
\chapters@set{20261125T010500}{9}% poincare-series-and-near-orthogonality
\chapters@set{20261124T223000}{10}% elliptic-curves-for-the-proof
\chapters@set{20261125T009500}{11}% fock-model-and-theta-intertwiners
\chapters@set{20261124T223100}{12}% complex-multiplication-and-inert-reduction
\chapters@set{20261125T008000}{13}% one-cm-determinant
\chapters@set{20261124T214500}{14}% supersingular-theta-divisibility
\chapters@set{20261124T214600}{15}% the-isotropic-plane
\chapters@set{20261125T013700}{16}% pairings-on-a-formal-group-of-height-two
\chapters@set{20261125T013500}{17}% hecke-nondegeneracy-on-the-tree
\chapters@set{20261126T030200}{18}% the-bound-is-attained
\chapters@set{20261125T009000}{19}% kronecker-limit-and-chowla-selberg
\chapters@set{20261125T002000}{20}% definite-theta-lifts
\chapters@set{20261125T015500}{21}% siegel-weil-in-the-smallest-case
\chapters@set{20261125T002100}{22}% definite-theta-lifts-and-norms
\chapters@set{20261125T016500}{23}% hecke-operators-and-theta-lifts
\chapters@set{20261125T016700}{24}% divisibility-along-the-chain
\chapters@set{20261125T004000}{25}% denominators-and-sha
\chapters@set{20261125T004100}{26}% tate-twists-and-sha
\chapters@set{20261126T020430}{27}% the-minus-part-by-gauss-sums
\chapters@set{20261125T014100}{28}% congruence-modules
\chapters@set{20261125T014500}{29}% eisenstein-congruences-and-selmer-classes
\chapters@set{20261125T014700}{30}% pseudocharacters-and-the-length-bound
\chapters@set{20261125T011000}{31}% galois-representations-black-box
\chapters@set{20261126T030300}{32}% the-data
\chapters@set{20261126T020210}{33}% bessel-poincare-series-and-their-gram-operators
\chapters@set{20261126T030400}{34}% the-seesaw-identity
\chapters@set{20261126T030500}{35}% divisibility-of-the-composite-lift
\chapters@set{20261126T030600}{36}% theorem-a-short-forms
\chapters@set{20261126T030700}{37}% the-determinant-of-a-chain
\chapters@set{20261126T030800}{38}% the-cost-of-one-lift
\chapters@set{20261126T030900}{39}% denominators-of-projections-to-lifts
\chapters@set{20261126T031000}{40}% twists-with-few-congruences
\chapters@set{20261126T031100}{41}% theorem-c-and-the-main-inequality
\providecommand\chapnumber[1]{%
  \@ifundefined{chapters@#1}%
    {\PackageError{chapters}{Sheet #1 is not listed in the index}%
      {List it in 20261125T012000--grh-workbook__guide.org\MessageBreak
       and run scripts/make-chapters.py.}}%
    {\edef\thechap{\@nameuse{chapters@#1}}}}
\makeatother
\chapnumber{<sheet ID>}
\makeatletter
\providecommand\chapters@set[2]{\@namedef{chapters@#1}{#2}}
\chapters@set{20261125T001000}{1}% zeros-to-ramp-sums
\chapters@set{20261125T001100}{2}% inert-primes-and-estimate-b
\chapters@set{20261125T005000}{3}% galois-fourier-and-averaging
\chapters@set{20261125T005500}{4}% columns-from-a-sphere
\chapters@set{20261125T000000}{5}% slope-method
\chapters@set{20261126T030100}{6}% a-twist-for-each-prime
\chapters@set{20261125T003000}{7}% the-chain-and-the-square-root
\chapters@set{20261125T010000}{8}% rank-lemma-and-zero-estimate
\chapters@set{20261125T010500}{9}% poincare-series-and-near-orthogonality
\chapters@set{20261124T223000}{10}% elliptic-curves-for-the-proof
\chapters@set{20261125T009500}{11}% fock-model-and-theta-intertwiners
\chapters@set{20261124T223100}{12}% complex-multiplication-and-inert-reduction
\chapters@set{20261125T008000}{13}% one-cm-determinant
\chapters@set{20261124T214500}{14}% supersingular-theta-divisibility
\chapters@set{20261124T214600}{15}% the-isotropic-plane
\chapters@set{20261125T013700}{16}% pairings-on-a-formal-group-of-height-two
\chapters@set{20261125T013500}{17}% hecke-nondegeneracy-on-the-tree
\chapters@set{20261126T030200}{18}% the-bound-is-attained
\chapters@set{20261125T009000}{19}% kronecker-limit-and-chowla-selberg
\chapters@set{20261125T002000}{20}% definite-theta-lifts
\chapters@set{20261125T015500}{21}% siegel-weil-in-the-smallest-case
\chapters@set{20261125T002100}{22}% definite-theta-lifts-and-norms
\chapters@set{20261125T016500}{23}% hecke-operators-and-theta-lifts
\chapters@set{20261125T016700}{24}% divisibility-along-the-chain
\chapters@set{20261125T004000}{25}% denominators-and-sha
\chapters@set{20261125T004100}{26}% tate-twists-and-sha
\chapters@set{20261126T020430}{27}% the-minus-part-by-gauss-sums
\chapters@set{20261125T014100}{28}% congruence-modules
\chapters@set{20261125T014500}{29}% eisenstein-congruences-and-selmer-classes
\chapters@set{20261125T014700}{30}% pseudocharacters-and-the-length-bound
\chapters@set{20261125T011000}{31}% galois-representations-black-box
\chapters@set{20261126T030300}{32}% the-data
\chapters@set{20261126T020210}{33}% bessel-poincare-series-and-their-gram-operators
\chapters@set{20261126T030400}{34}% the-seesaw-identity
\chapters@set{20261126T030500}{35}% divisibility-of-the-composite-lift
\chapters@set{20261126T030600}{36}% theorem-a-short-forms
\chapters@set{20261126T030700}{37}% the-determinant-of-a-chain
\chapters@set{20261126T030800}{38}% the-cost-of-one-lift
\chapters@set{20261126T030900}{39}% denominators-of-projections-to-lifts
\chapters@set{20261126T031000}{40}% twists-with-few-congruences
\chapters@set{20261126T031100}{41}% theorem-c-and-the-main-inequality
\providecommand\chapnumber[1]{%
  \@ifundefined{chapters@#1}%
    {\PackageError{chapters}{Sheet #1 is not listed in the index}%
      {List it in 20261125T012000--grh-workbook__guide.org\MessageBreak
       and run scripts/make-chapters.py.}}%
    {\edef\thechap{\@nameuse{chapters@#1}}}}
\makeatother

\ifcsname chapters@20261126T081200\endcsname
  \chapnumber{20261126T081200}
\else
  \renewcommand{\thechap}{?}
\fi
\renewcommand{\thesection}{\thechap.\arabic{section}}
\theoremstyle{definition}
\newtheorem{problem}{Problem}[section]
\newtheorem{definition}{Definition}[section]
\theoremstyle{plain}
\newtheorem{theorem}{Theorem}[section]
\theoremstyle{remark}
\newtheorem*{leanremark}{In the Lean formalization}

\newcommand{\Z}{\mathbb{Z}}
\newcommand{\Q}{\mathbb{Q}}
\newcommand{\R}{\mathbb{R}}
\newcommand{\C}{\mathbb{C}}
\newcommand{\F}{\mathbb{F}}
\newcommand{\Ql}{\bar{\mathbb{Q}}_\ell}
\newcommand{\Zl}{\bar{\mathbb{Z}}_\ell}
\newcommand{\ii}{\mathrm{i}}
\newcommand{\Oc}{\mathcal{O}}
\newcommand{\Mcal}{\mathcal{M}}
\newcommand{\Nc}{\mathcal{N}}
\newcommand{\St}{\mathrm{St}}
\newcommand{\bd}{\mathrm{bd}}
\newcommand{\ms}{\mathrm{ms}}
\newcommand{\alg}{\mathrm{alg}}
\newcommand{\Fr}{\mathrm{Fr}}
\newcommand{\Cl}{\mathrm{Cl}}
\newcommand{\Nm}{\mathrm{Nm}}
\newcommand{\val}{\operatorname{val}}
\newcommand{\tl}[1]{\texttt{#1}}

\setlength{\emergencystretch}{3em}

\title{A measure for Hecke values}
\author{}
\date{November 26, 2026}

\begin{document}
\maketitle
\refstepcounter{chap}\label{chap}

\paragraph{Prerequisites.}
Weierstrass functions and local torsion \cite[Thm.~VII.3.4, Prop.~VII.4.1]{Si1}; cyclic norms \cite{La}.  Chapter~\ref{ec-chap}: Probs.~\ref{ec-pr:wp-laurent}, \ref{ec-pr:square}, \ref{ec-pr:varpi}, \ref{ec-pr:hurwitz}.  Chapter~\ref{da-chap}: the data \(D\) (\(E\) with \(f\) infinite places, \(r\), \(m\), \(M = p_0p_1\), \(N\)).  Chapter~\ref{ob-chap}: \(\val\), \(\val_+\), \(\Oc^{[q]}\); \(\Mcal\), \(\Nc\) (\S\ref{ob-sec:chosen}), gauge (Def.~\ref{ob-def:gauge}), tame characters and twists (\S\ref{ob-sec:twists}), \(\upsilon\) (\S\ref{ob-sec:characters}), critical pairs and ratios (Def.~\ref{ob-def:pair}).  Chapter~\ref{dp-chap}: \(G_E := \mathrm{Gal}(\bar{\Q}/E)\), \(s_w\), \(\mathrm{Art}\), \(\val_\ell\) (\S\ref{dp-sec:statement}), \(\mathrm{tw}_\ell\) (Def.~\ref{dp-def:clearing-bound}), \(\hat\xi\) and auxiliary twists (Def.~\ref{dp-def:aux-twist}), Defs.~\ref{dp-def:main-summand-like}, \ref{dp-def:selmer-bound}.  Chapter~\ref{l3-chap}: \(\phi_{P,\xi}\), \(d_P\), \(\alg_P\), Thm.~\ref{l3-thm:pair}.  Chapter~\ref{c3-chap}: \(\mathrm{ex}_h\) (Def.~\ref{c3-def:exchange-factor}).  Chapter~\ref{w2-chap}: characters for \(Q\), \(\Cl_Q\), Stickelberger systems (Defs.~\ref{w2-def:chi}, \ref{w2-def:system}), Probs.~\ref{w2-pr:euler-line}, \ref{w2-pr:ratio-from-system}, Thm.~\ref{w2-thm:polynomials}; Chapter~\ref{w3-chap}: Thm.~\ref{w3-thm:tied}; Chapter~\ref{tw-chap}: Thm.~\ref{tw-thm:twisting}.

\medskip
In the twisting theorem (Theorem~\ref{tw-thm:twisting} of Chapter~\ref{tw-chap}), fix the data \(D\) of Chapter~\ref{da-chap}, a prime \(\ell\), \(\iota \colon \C \simeq \Ql\), and a finite family \(R_i\) of main-summand-like \(\ell\)-adic representations of \(G_E\) (Definition~\ref{dp-def:main-summand-like}).  Write \(\mathrm{tw}_\ell := (4f + 3)r\,\val_\ell(\Nc) + 2^{-\ell}\), as in Definition~\ref{dp-def:clearing-bound}.  The theorem gives one \(m\)-fold twist \(\vartheta\) such that, for each \(1 \le j \le m\), \(\vartheta^j\) is auxiliary with ramification at most \(\mathrm{tw}_\ell\) (Definition~\ref{dp-def:aux-twist}), every twisted representation \(\widehat{\vartheta^j}R_i\) has Selmer bound \(\mathrm{tw}_\ell\) (Definition~\ref{dp-def:selmer-bound}), and \(\val_+(\iota(R_P(\vartheta^j, \epsilon))) \le \mathrm{tw}_\ell\) for every ramified critical pair \(P = (\nu, z)\) and each sign \(\epsilon = \pm 1\).  The proof uses a Stickelberger system for each ramified pair and sign, the polynomial input of Theorem~\ref{w2-thm:polynomials}, and the simultaneous choice of roots of unity in Chapter~\ref{tw-chap}.

How do twisted partial Hecke \(L\)-values give a measure (Definition~\ref{def:measure}) whose coefficients at modulus \(Q\) are integral away from \(\Nm Q\)?  How does this measure give elements \(\St_Q \in \Zl[\Cl_Q]\) obeying Euler-factor insertion, and where does the construction use ramification at the chosen place \(w_0\)?  This chapter works out unramified and ramified models over \(\Q(\ii)\), defines the measure of a critical pair, and states Theorem~\ref{thm:pair-system}.

Throughout, \(\ell\) is a prime, \(\val\) is the valuation of \(\Ql\) with \(\val(\ell) = 1\) and \(\val(0) = \infty\), \(\Zl\) is its valuation ring, \(\val_+ := \max(\val, 0)\), \(\lVert y\rVert := \ell^{-\val(y)}\), and \(\val_\ell\) is the \(\ell\)-adic valuation on \(\Q\); inner products \(\langle\cdot, \cdot\rangle\) are conjugate-linear in the first variable.

\section{Unramified and ramified models over \texorpdfstring{\(\Q(\ii)\)}{Q(i)}}\label{sec:models}

Since \(\Q(\ii)\) is not a field of the data of Chapter~\ref{da-chap}, the constructions in this section are models.  Write \(\ii := \sqrt{-1}\), \(\mathrm{N}(z) := z\bar z\) on \(\Q(\ii)\), and \(\Omega := 2\int_0^1 (1 - t^4)^{-1/2}\,dt\) for the lemniscate period (\(\varpi\) in Chapter~\ref{ec-chap}, distinct from \(\Omega\) in Chapter~\ref{da-chap}), so \(\sum_{0 \ne \lambda \in \Z[\ii]} \lambda^{-4} = \Omega^4/15\) by Hurwitz's formula (Problem~\ref{ec-pr:hurwitz}; \emph{status}: classical).  For a prime \(q \equiv 1 \pmod 4\), fix \(\pi \in \Z[\ii]\) with \(q = \pi\bar\pi\), identify \(\Z[\ii]/(\pi) \cong \F_q\), and put \(\Cl_q := \F_q^\times / \{\pm 1\}\), cyclic of order \((q-1)/2\); for \(\lambda \in \Z[\ii]\) with \((\lambda, q) = 1\), let \([\lambda] \in \Cl_q\) be the coset of \(\lambda/\bar\lambda \bmod \pi\).  For a finite set \(Q\) of places \(w = (\pi_w)\) above distinct primes \(q_w \equiv 1 \pmod 4\), put \(\Nm Q := \prod_{w \in Q} q_w\), \([\lambda]_Q := ([\lambda]_{q_w})_{w \in Q}\), and \(\St_Q := (15/\Omega^4) \sum_{\substack{0 \ne \lambda \in \Z[\ii] \\ (\lambda, \Nm Q) = 1}} \lambda^{-4}\,[\lambda]_Q \in \C\bigl[\prod_{w \in Q} \Cl_{q_w}\bigr]\), so \(1(\St_\varnothing) = 1\) at the trivial character \(1\); for \(Q = \{(\pi)\}\) write \(\St_q = \sum_{u \in \Cl_q} c_u\,[u]\) with \(c_u := (15/\Omega^4) \sum_{(\lambda, q) = 1,\, [\lambda] = u} \lambda^{-4} \in \C\).

Recall (Definition~\ref{w2-def:system}) that a Stickelberger system for \((\Sigma_0, g_0)\) over \(E\) has a continuous character \(\alpha \colon G_E \to \Ql^\times\) unramified outside \(\Sigma_0\) with \(\alpha(g_0) \ne 1\), a finite list \(L\) of integer exponents, and an element \(\St_Q \in \Zl[\Cl_Q]\) for each finite set \(Q\) of finite places, with \(1(\St_\varnothing) \ne 0\), obeying Euler-factor insertion when \(v \notin Q\) and \(Q \cup \{v\}\) is admissible for \(\Sigma_0\).  The first problem checks that \([\lambda]\) and \(c_u\) are well defined, that \(\chi(\St_q)\) is real, and how \(\Cl_q\) maps onto \(\Cl_{\{w\}}\).

\begin{problem}\label{pr:clq}
\textup{(a)} For \(\lambda \in \Z[\ii]\) prime to \(q\), \([\lambda] \in \Cl_q\) depends only on \(\lambda\,\Z[\ii]^\times\) and satisfies \([\bar\lambda] = [\lambda]^{-1}\); each series \(c_u\) converges absolutely with \(\bar c_u = c_{u^{-1}}\), so \(\chi(\St_q) := \sum_{u \in \Cl_q} c_u\,\chi(u) \in \R\) for every character \(\chi\) of \(\Cl_q\).
\textup{(b)} For \(E = \Q(\ii)\) and \(w = (\pi)\), reduction modulo \(\pi\) on norm-one finite ideles that are units everywhere identifies \(\Cl_{\{w\}}\) (\S\ref{w2-sec:systems}) with \(\Cl_q / \langle \ii \bmod \pi\rangle\); hence every odd-order character of \(\Cl_q\) factors through \(\Cl_{\{w\}}\).
\end{problem}

The next problem proves the insertion law for \(\St_Q\) and evaluates \(1(\St_q)\).

\begin{problem}\label{pr:insert}
For \(Q\) as above, \(v = (\pi_v)\) above a prime \(q_v \equiv 1 \pmod 4\) not dividing \(\Nm Q\), \(a_v := (\bar\pi_v/\pi_v)^2\), and any character \(\chi\) of \(\prod_{w \in Q} \Cl_{q_w}\) pulled back to \(Q \cup \{v\}\),
\[
  \chi(\St_{Q \cup \{v\}}) = \chi(\St_Q) \prod_{\epsilon = \pm 1} \bigl(1 - (a_v\,\chi([\pi_v]_Q))^{\epsilon}\, q_v^{-2}\bigr),
\]
of the shape of Definition~\ref{w2-def:system} with \(L = [2]\).  In particular, \(1(\St_q) = (1 - \pi^{-4})(1 - \bar\pi^{-4}) = \mathrm{N}(\pi^4 - 1)/q^4\), equal to \(128/125\) and \(28800/28561\) for \(q = 5, 13\).
\end{problem}

Let \(\wp\) be the Weierstrass function of \(\Omega\Z[\ii]\), uniformizing \(y^2 = 4x^3 - 4x\) (Problems~\ref{ec-pr:wp-laurent}, \ref{ec-pr:varpi}); \((X, Y) = (x, y/2)\) gives \(Y^2 = X^3 - X\), of discriminant \(64\).  For a lift \(\tilde a \in \Z[\ii]\) of \(a \in (\Z[\ii]/q)^\times\), both \([a] := [\tilde a] \in \Cl_q\) and \(x(P_a) := \wp(\Omega\tilde a/q)\) are independent of \(\tilde a\).  Over a local field of characteristic zero, a nonzero point \((X, Y)\) of odd prime order \(q\) on an integral Weierstrass equation has \(X\) integral when \(q\) is prime to the residue characteristic, and \(\varpi^{2\lfloor e/(q-1)\rfloor} X\) integral for a uniformizer \(\varpi\) of the local field, where \(e\) is the field's absolute ramification index, when \(q\) equals the residue characteristic \cite[Thm.~VII.3.4]{Si1} (\emph{status}: classical).  The next problem expresses \(c_u\) via \(x(P_a)\) and deduces its local integrality.

\begin{problem}\label{pr:weierstrass}
\textup{(a)} For every \(u \in \Cl_q\), \(c_u = (5/q^4) \sum_{[a] = u} \bigl(3x(P_a)^2 - 1\bigr)\) over \(a \in (\Z[\ii]/q)^\times\).
\textup{(b)} Hence \(q^4 c_u / 5\) is integral away from \(q\), \(c_u \in \Oc^{[q]}\), and \(\iota(c_u) \in \Zl\) for every \(\ell \ne q\) and \(\iota \colon \C \simeq \Ql\); moreover \(q\,x(P_a)\) and \(q^6 c_u\) are algebraic integers.
\end{problem}

The next problem shows that the multiset of coefficients \(c_u\) is Galois-invariant over \(\Q\) and finds its polynomial for small \(q\).

\begin{problem}\label{pr:poly}
\textup{(a)} Every \(\sigma \in \mathrm{Gal}(\bar{\Q}/\Q)\) permutes \((c_u)_{u \in \Cl_q}\), so \(\prod_{u \in \Cl_q}(X - c_u) \in \Q[X]\) and \(q^{6k}\) clears the coefficient of \(X^{(q-1)/2 - k}\) into \(\Z\).

\textup{(b)} For \(q = 5\) (\(\pi = 2 + \ii\)), \(\prod_{u \in \Cl_5}(X - c_u) = X^2 - \frac{128}{125}X - \frac{7424}{15625}\) and \(c_{\pm 1} - c_{\pm 2} = 96\sqrt{5}/125\).

\textup{(c)} By machine, for \(q = 13, 17\) the coefficient of \(X^{(q-1)/2 - k}\) has reduced denominator \(q^{4k}\) (\(0 \le k \le (q-1)/2\)).
\end{problem}

The next problem computes the \(3\)-adic valuations of \(\chi(\St_{13})\) at orders \(1\) and \(3\).

\begin{problem}\label{pr:unram-vals}
For \(d \ge 3\) dividing \((q-1)/2\), the product \(P_d(q) := \prod_{\chi} \chi(\St_q)\) over the \(\varphi(d)\) characters \(\chi\) of \(\Cl_q\) of exact order \(d\) lies in \(\Q\) with \(q^{6\varphi(d)} P_d(q) \in \Z\).  By machine, \(P_3(13) = 67737600/13^7 = 2^{11}\cdot 3^3\cdot 5^2\cdot 7^2/13^7\); under every \(\iota \colon \C \simeq \bar{\Q}_3\), \(\val(\iota(\chi(\St_{13})))\) is \(2\) at \(\chi = 1\) and \(3/2\) at each character of order \(3\).
\end{problem}

The next problem shows why \(1(\St_q)\) has \(3\)-adic valuation at least \(2\) for \(q \equiv 1 \pmod{12}\).

\begin{problem}\label{pr:price}
For every prime \(q \equiv 1 \pmod{12}\) with \(q = \pi\bar\pi\) in \(\Z[\ii]\), \(\pi^4 \equiv \mathrm{N}(\pi) \equiv 1 \pmod 3\) in \(\Z[\ii]\); hence \(9 \mid \mathrm{N}(\pi^4 - 1)\) and \(\val_3(1(\St_q)) \ge 2\).
\end{problem}

A split prime \(q\) admitting a nontrivial tame residue character of \(3\)-power order satisfies \(q \equiv 1 \pmod{12}\).  Problem~\ref{pr:price} shows that inserting such a prime in the unramified model increases the valuation at the trivial character by at least \(2\).  The Hecke character supplying \(a_v\) in Problem~\ref{pr:insert} is unramified at every finite place, so its \(\ell\)-adic avatar is trivial on inertia at every finite place away from \(\ell\).  We now introduce a model with a nontrivial tame part at \(w_0 := (2 + \ii)\) above \(5\), and compute \(3\)-adic unit ratios for the example at \(q = 13\).

Identify \(\Z[\ii]/w_0 \cong \F_5\) via \(\ii \mapsto 3\), and write \((x \mid w_0) \in \{\pm 1\}\) for the quadratic character of \(\F_5^\times\).  For \(Q\) as above with \(5 \nmid \Nm Q\), put \(\St^\nu_Q := (15/\Omega^4) \sum_{(\lambda, 5\,\Nm Q) = 1} (\lambda/\bar\lambda \mid w_0)\,\lambda^{-4}\,[\lambda]_Q\) (written \(\St^\nu_q\) for \(Q = \{(\pi)\}\), \(q = \pi\bar\pi \ne 5\)).  Let \(\psi_1\) be the algebraic Hecke character with ideal values \(\psi_1((\lambda)) := (\lambda/\bar\lambda \mid w_0)(\lambda/\bar\lambda)^{-2}\) on nonzero ideals prime to \(5\).  For \(\ell \ne 5\) and \(\iota \colon \C \simeq \Ql\), let \(\alpha\) be the \(\ell\)-adic avatar of \(\psi_1^{-1}\) through \(\iota\).  At a tame inertia generator \(g_0\) at \(w_0\) it has \(\alpha(g_0) = -1\).  For \(L = [2]\), put \(\bd(\St^\nu) := \val(\iota(1(\St^\nu_\varnothing))) + 2\val(\alpha(g_0) - 1)\).  The next problem evaluates \(1(\St^\nu_\varnothing)\), \(\bd(\St^\nu)\) and the untwisted factor at \(q = 13\).

\begin{problem}\label{pr:ram-empty}
\textup{(a)} From \(\St_5\), \(1(\St^\nu_\varnothing) = c_{\pm 1}(5) - c_{\pm 2}(5) = 96\sqrt{5}/125\); for \(\ell = 3\) and every \(\iota \colon \C \simeq \bar{\Q}_3\), \(\val(\iota(1(\St^\nu_\varnothing))) = \bd(\St^\nu) = 1\).
\textup{(b)} For a prime \(q = \pi\bar\pi \equiv 1 \pmod 4\), \(q \ne 5\), and \(s := (\pi/\bar\pi \mid w_0)\), \(1(\St^\nu_q) = 1(\St^\nu_\varnothing)(1 - s\pi^{-4})(1 - s\bar\pi^{-4})\); at \(\pi = 3 + 2\ii\) (\(q = 13\)), \(s = -1\) and \((1 - s\pi^{-4})(1 - s\bar\pi^{-4}) = 28324/13^4 = 2^2\cdot 73\cdot 97/13^4\), a \(3\)-adic unit.
\end{problem}

The next problem identifies \(\chi(\St^\nu_q)/1(\St^\nu_\varnothing)\) with an \(L\)-value ratio and checks that the order-\(3\) ratios at \(q = 13\) are \(3\)-adic units.

\begin{problem}\label{pr:ram-ratios}
\textup{(a)} For a prime \(q = \pi\bar\pi \equiv 1 \pmod 4\), \(q \ne 5\), and a character \(\chi\) of \(\Cl_q\), define \(\psi_\chi((\lambda)) := \psi_1((\lambda))\,\chi([\lambda])\) on nonzero ideals prime to \(5q\), and \(L^S(2, \psi) := \prod_{\mathfrak p \notin S \cup \bar S}(1 - \psi(\mathfrak p)\,\mathrm{N}(\mathfrak p)^{-2})^{-1}\) over nonzero prime ideals \(\mathfrak p\) of \(\Z[\ii]\), with \(\mathrm{N}(\mathfrak p) := \#(\Z[\ii]/\mathfrak p)\).  Then \(\chi(\St^\nu_q)/1(\St^\nu_\varnothing) > 0\) equals \(L^{\{w_0, (\pi)\}}(2, \psi_\chi)/L^{\{w_0\}}(2, \psi_1)\).

\textup{(b)} On a norm-one idele \(u\) that is a unit at every finite place, the formula \(\nu(u) := u_\infty^{-2}(u_{w_0} \bmod w_0 \mid w_0)\), where \(u_\infty\) is the infinite component, defines a character on the norm-one idele classes.  Let \(P := (\nu, 2)\) be the model pair.  If \(\chi\) factors through \(\Cl_{\{(\pi)\}}\) in Problem~\ref{pr:clq}\textup{(b)}, this ratio is the model of \(R_P(\xi, 1)\) for the tame character \(\xi\) of modulus \(\{(\pi)\}\) with character \(\chi\) on \(\Cl_{\{(\pi)\}}\) (tame part \(\chi\) at \((\pi)\)).

\textup{(c)} By machine, the product of \(\chi(\St^\nu_{13})/1(\St^\nu_\varnothing)\) over the two characters \(\chi\) of order \(3\) is \(62939569/13^7 = 7^2\cdot 11\cdot 23\cdot 5077/13^7\), and both order-\(3\) ratios are \(3\)-adic units.
\end{problem}

\section{Measures and Stickelberger systems}\label{sec:measures}

We return to the data \(D\) of Chapters~\ref{da-chap} and~\ref{ob-chap}: the CM field \(E\) with infinite places \(\Phi\), \(M = p_0p_1\), \(N\), \(\Mcal\), \(\Nc\), tame characters \(\tau\) of modulus \(Q_\tau\) and norm \(\Nm Q_\tau\) (\S\ref{ob-sec:twists}; \(\tau\) is a tame character below), twists \(\xi\), and critical pairs \(P = (\nu, z)\) of negative weight \(\mathrm{wt}_P \colon \Phi \to \Z_{\ge 0}\) (ramified when the tame part of \(\nu\) is nontrivial precisely on \(\Mcal\)), with \(L^S(z, \mu)\) and \(R_P(\xi, \epsilon) = \lvert L^{\Mcal \cup Q_\xi}(z, (\nu\xi^{\epsilon})^{-1})\rvert / \lvert L^{\Mcal}(z, \nu^{-1})\rvert\) (Definition~\ref{ob-def:pair}).  Put \((L_z, n_z) := ([0, 1], 2)\) if \(z = 1\) and \(([z], 1)\) if \(z \ge 2\); for \(v \notin \Mcal \cup \bar\Mcal \cup Q_\tau \cup \bar Q_\tau\), let \(c_v(\nu\tau)\) be the value of \(x \mapsto (\nu\tau)(x/\bar x)\) at a uniformizer at \(v\).  The following definition states the axioms of a measure.

\begin{definition}\label{def:measure}
Let \(P = (\nu, z)\) be a critical pair.  A \emph{measure} of \(P\) is a function \(\ms_P\) from tame characters to \(\C\) such that:
\begin{enumerate}
\item \(\ms_P(1) \ne 0\), where \(1\) is the trivial tame character (\(Q_1 = \varnothing\));
\item for every modulus \(Q\) of a tame character, there are finitely many norm-one ideles \(u_i\) of \(E\) and \(c_i \in \Oc^{[\Nm Q]}\) such that \(\ms_P(\tau) = \sum_i c_i\,\tau(u_i)\) whenever \(Q_\tau = Q\);
\item if \(\tau, \tau'\) have the same underlying character and tame part and \(Q_{\tau'} = Q_\tau \cup \{v\}\) for \(v \notin Q_\tau\), then
\[
  \ms_P(\tau') = \ms_P(\tau) \prod_{e \in L_z} \prod_{\epsilon = \pm 1} \bigl(1 - c_v(\nu\tau)^{\epsilon}\,(\#k_v)^{-e}\bigr) ;
\]
\item for every twist \(\xi\), \(\ms_P(\xi) = \ms_P(1)\,(\Nm Q_\xi)^{n_z - 1}\,R_P(\xi, 1)^{n_z}\).
\end{enumerate}
\end{definition}

Fix \(\iota \colon \C \simeq \Ql\), let \(\Sigma_0\) be the finite places of \(E\) dividing \(\ell MN\), let \(w_0 \in \Mcal\) not divide \(\ell\), and let \(g_0 = s_{w_0}\) be its tame generator (\S\ref{dp-sec:statement}).  A character \(\chi\) for \(Q\) has \(\chi([y/\bar y]) = \chi(\mathrm{Art}\,y)\) on \(\Cl_Q\) (\S\ref{w2-sec:systems}), and \(\hat\xi(\mathrm{Art}\,y) = \iota(\xi(y/\bar y))^{-1}\) (Definition~\ref{dp-def:aux-twist}); for \(\epsilon = \pm 1\), \(\xi^{\epsilon}\) denotes the tame character of modulus \(Q_\xi\) with character \(\xi^{\epsilon}\) and tame part \((\psi_w^{\epsilon})\).  The next problem bounds \(\bd(\St) = \val(1(\St_\varnothing)) + 2\,\#L\,\val(\alpha(g_0) - 1)\) at a root of unity \(\alpha(g_0) \ne 1\).

\begin{problem}\label{pr:root-val}
For every \(n \ge 1\), \(\prod_{\zeta^n = 1,\, \zeta \ne 1} (1 - \zeta) = n\) in \(\Ql\), so every \(u \ne 1\) in \(\Ql\) with \(u^n = 1\) satisfies \(\val(u - 1) \le \val_\ell(n)\).  Hence any Stickelberger system \(\St\) for \((\Sigma_0, g_0)\) with \(\alpha(g_0)^{\#k_{w_0} - 1} = 1\) and \(\#L \le 2\) satisfies \(\bd(\St) \le \val(1(\St_\varnothing)) + 4\val_\ell(\#k_{w_0} - 1)\).
\end{problem}

The next problem passes from a measure at a fixed modulus prime to \(\ell\) to \(\Zl[\Cl_Q]\).

\begin{problem}\label{pr:measure-to-st}
Every finite idele \(u\) of \(E\) with \(u\bar u = 1\) equals \(y/\bar y\) for a finite idele \(y\).  Let \(\ms_P\) be a measure of \(P\), \(\epsilon = \pm 1\), and \(Q\) the modulus of a tame character with \(\ell \nmid \Nm Q\); write \(\ms_P(\tau) = \sum_i c_i\,\tau(u_i)\) on modulus \(Q\), with finite part \(y_i/\bar y_i\) for each \(u_i\).  Then \(\St_Q := \sum_i \iota(c_i)\,\bigl[(y_i/\bar y_i)^{-\epsilon}\bigr] \in \Zl[\Cl_Q]\), and \(\hat\xi(\St_Q) = \iota(\ms_P(\xi^{\epsilon}))\) for every tame character \(\xi\) of modulus \(Q\) whose \(\hat\xi\) is a character for \(Q\).
\end{problem}

\section{Norms of lifts, a period at \texorpdfstring{\(1\)}{1}, and the systems of the ratios}\label{sec:systems}

Let \(P = (\nu, z)\) be a critical pair with \(z \ge 2\), with lift \(\phi_{P,\tau}\) in rank \(2z\), denominator \(d_P\), and \(\alg_P(\tau) = \Nc\,d_P^2\,\lVert\phi_{P,\tau}\rVert^2/(\Nm Q_\tau)^{2z-1} \in \Oc^{[\Nm Q_\tau]}\) satisfying \(R_P(\tau, 1) = \alg_P(\tau)/\alg_P(1)\) (Theorem~\ref{l3-thm:pair}).  For tame characters \(\tau\) of a fixed modulus \(Q\), the proof writes the value of the lift at each finite adelic point \(h\) of the target as \(\phi_{P,\tau}(h) = \sum_i \tau(\det g_i)^{-1} K_i(h)\).  Here the \(g_i\) are adelic points of \(\mathrm U(1)\) with infinite component \(1\), and \(K_i(h)\) is the scaled theta kernel at \((g_i, h)\) for the trivial character regarded as a tame character of modulus \(Q\).  The points and kernel functions are fixed once \(P\) and \(Q\) are fixed.  At every \(h\) and every ring homomorphism \(\eta \colon \C \to \C_\ell\) for a prime \(\ell \nmid \Nm Q\), the value \(d_P K_i(h)\) has gauge at most \(1\).  The next problem gives the pointwise inner product identity.

\begin{problem}\label{pr:norm-measure}
For \(K_1, \dots, K_n\) in a complex inner product space \(V\), \(g_1, \dots, g_n\) in the adelic points of a unitary group for \(E/E^+\), where \(E^+\) is the maximal real subfield of \(E\), and a unitary character \(\tau\) of the norm-one ideles of \(E\), \(T(\tau) := \sum_{i=1}^n \tau(\det g_i)^{-1} K_i\) has \(\lVert T(\tau)\rVert^2 = \sum_{i, j = 1}^n \langle K_j, K_i\rangle\,\tau(\det g_j / \det g_i)\).
\end{problem}

At fixed representatives of the classes for the level of the lift, let \(B_{ji}\) be the class sum of the Gram pairings \(\langle K_j(h), K_i(h)\rangle\), weighted by the inverse orders of the class stabilizers.  The proof uses the coefficients \(\Nc\,d_P^2 B_{ji}/(\Nm Q)^{2z-1}\), which are integral away from \(\Nm Q\), for clause~(2) of Definition~\ref{def:measure}.  Together with Theorem~\ref{l3-thm:pair} and the deletion of the Euler factors at \(v\) and \(\bar v\) when the same tame character is regarded at modulus \(Q_\tau \cup \{v\}\), this gives that \(\ms_P := \alg_P\) is a measure of \(P\) for \(z \ge 2\).  \emph{Status}: read from the proof; not proved here.

When \(z = 1\), rank \(2z = 2\) falls outside Theorem~\ref{l3-thm:pair}.  For a ramified critical pair \(P = (\nu, 1)\), let \(\Lambda(\tau)\) be the determinant theta period at the standard level on the definite unitary group of rank \(2\), for the character \(\nu\tau\upsilon\), the determinant exponents \(\mathrm{wt}_P(\sigma) - 1\) (\(\sigma \in \Phi\)), and the product of the tame parts of \(\nu\) and \(\tau\).  Its datum has modulus \(\prod_{w \in \Mcal} w^2 \prod_{w \in Q_\tau} w\).  The proof chooses a nonzero scalar \(N_P \in \C\) independent of \(\tau\) and its modulus, with which \(\ms_P(\tau) := \lvert\Nc N_P \Lambda(\tau)\rvert^2\) is a measure with \(L_1 = [0, 1]\) and \(n_1 = 2\).  There is a real constant \(C_P \ne 0\) such that, for every twist \(\xi\), \(\Lambda(\xi) \ne 0\) and \(C_P\lvert\Lambda(\xi)\rvert^2 = \mathrm{ex}_1(\Mcal \cup Q_\xi, (\nu\xi)^{-1})\), with \(\mathrm{ex}_1\) as in Definition~\ref{c3-def:exchange-factor}.  If \(\tau'\) has the same underlying character and tame part as \(\tau\) and \(Q_{\tau'} = Q_\tau \cup \{v\}\) for \(v \notin Q_\tau\), then \(\Lambda(\tau') = (1 - X)(1 - qX)\Lambda(\tau)\), where \(q = \#k_v\) and \(X = c_v(\nu\tau)^{-1}/q\).  \emph{Status}: read from the proof; not proved here (for \(p\)-adic CM measures, see \cite{Ka,Hi}).  The next problem checks the squared-modulus Euler identity for \(L_1 = [0, 1]\).

\begin{problem}\label{pr:z-one}
For \(q > 0\), \(a \in \C\) with \(\lvert a\rvert = 1\) and \(X := a^{-1}/q\), \(\lvert (1 - X)(1 - qX)\rvert^2 = \prod_{e \in \{0, 1\}} \prod_{\epsilon = \pm 1} (1 - a^{\epsilon} q^{-e})\); in \(\Ql\) for \(\ell = 3\), \(q = 19\) and \(a = \zeta_9\) of exact order \(9\), the \(e = 1\) pair and the full product have valuations \(1/3\) and \(2/3\).
\end{problem}

For a ramified pair \(P = (\nu, z)\) and \(\epsilon = \pm 1\), the proof takes \(\alpha \colon G_E \to \Ql^\times\) to be the \(\ell\)-adic avatar via \(\iota\) of the base change of \(\nu^{-\epsilon}\): \(\alpha\) is continuous and unramified outside \(\Sigma_0\), satisfies \(\alpha(g_0) \ne 1\) and \(\alpha(g_0)^{\#k_{w_0} - 1} = 1\) since \(P\) is ramified, and has \(\alpha(\Fr_v)\,\hat\tau(\Fr_v) = \iota(c_v(\nu\tau^{\epsilon}))^{\epsilon}\) at arithmetic Frobenius \(\Fr_v\) for \(v \notin \Sigma_0 \cup Q_\tau \cup \bar Q_\tau\); moreover every character for an admissible \(Q\) is \(\hat\tau\) for some \(\tau\) with \(Q_\tau = Q\).  \emph{Status}: read from the proof; not proved here.  The following theorem supplies the Stickelberger systems for Chapters~\ref{w2-chap} and~\ref{tw-chap}.

\begin{theorem}[the systems of the ratios]\label{thm:pair-system}
Use the data and notation of Chapter~\ref{ob-chap}, and let \(\iota \colon \C \simeq \Ql\).  Let \(\Sigma_0\) be the finite places of \(E\) dividing \(\ell MN\), \(w_0\) a chosen place not above \(\ell\), and \(g_0 = s_{w_0}\) its tame generator at the fixed place above \(w_0\).  For every ramified critical pair \(P = (\nu, z)\) and sign \(\epsilon = \pm 1\), put \(n_z := 2\) if \(z = 1\) and \(n_z := 1\) otherwise; there is a Stickelberger system \(\St\) for \((\Sigma_0, g_0)\) with \(\bd(\St) \le \val(1(\St_\varnothing)) + 4\val_\ell(\#k_{w_0} - 1)\) such that, for every twist \(\xi\) whose set \(Q_\xi\) is admissible for \(\Sigma_0\) and whose inverse avatar \(\hat\xi\) is a character for \(Q_\xi\),
\[
  \bigl\lVert \hat\xi(\St_{Q_\xi}) \bigr\rVert = \bigl\lVert 1(\St_\varnothing) \bigr\rVert \cdot \bigl\lVert \iota\bigl(R_P(\xi, \epsilon)\bigr) \bigr\rVert^{n_z} .
\]
\end{theorem}

\emph{Status}: read from the proof; not proved here.  Problems~\ref{pr:clq}--\ref{pr:ram-ratios} establish in the \(\Q(\ii)\) models the insertion law for \(\St_Q\) and its ramified analogue at the trivial character, torsion integrality, the \(3\)-adic loss for \(q \equiv 1 \pmod{12}\), and the \(3\)-adic unit ratios in the example at \(q = 13\) of the sum ramified at \(w_0 = (2 + \ii)\); Problems~\ref{pr:norm-measure} and~\ref{pr:z-one} give the inner-product sum and the \(z = 1\) Euler identity.  Problems~\ref{pr:root-val} and~\ref{pr:measure-to-st} give the bound on \(\bd(\St)\) and the evaluation of the group-ring element at a fixed modulus.  In the proof the elements are taken to be zero at nonadmissible sets, and the norm identity uses clause~(4) of Definition~\ref{def:measure}.  Taken as given: Hurwitz's formula, the local bounds on torsion coordinates stated above \cite[Thm.~VII.3.4]{Si1}, the unramifiedness of torsion of order prime to the residue characteristic at places of good reduction \cite[Prop.~VII.4.1]{Si1}, cyclic norms \cite{La}, Theorem~\ref{l3-thm:pair}, and, from the proof, \(\ms_P\), the avatar facts, and Theorem~\ref{thm:pair-system}, which with Problem~\ref{w2-pr:ratio-from-system} and Theorems~\ref{w2-thm:polynomials} and~\ref{w3-thm:tied} supplies the Stickelberger input to Theorem~\ref{tw-thm:twisting}.

\begin{leanremark}
\raggedright In the Lean proof \cite{GL} (\tl{grh/src/}, namespace \tl{New.} omitted): Theorem~\ref{thm:pair-system} is \tl{Selmer.exists\_stickelbergerSystem\_}\(\beta\)\tl{\_le\_and\_norm\_eq}.  The lemma \tl{Selmer.ellVal\_le\_of\_norm\_eq\_mul\_pow} states that if \(x_0 \in \Zl\) is nonzero, \(i \ge 1\) is an integer, \(c \in [0,\infty]\), \(\lVert x\rVert = \lVert x_0\rVert\lVert y\rVert^i\), and \(\val_+(x) \le \val_+(x_0) + c\), then \(\val_+(y) \le c\).  The corollary \tl{Selmer.exists\_stickelbergerSystem\_ellVal\_ratio\_le} applies this lemma and the bound on \(\bd(\St)\) to obtain the implication in Problem~\ref{w2-pr:ratio-from-system}, with an allowance in \([0,\infty]\).  Definition~\ref{def:measure} specializes \tl{Data.CriticalPair.IsTwistMeasure} (which takes a finite set of nonnegative integer exponents and a nonnegative integer \(m\), with clause~(2) given by \tl{Data.TameChar.IsIntegralMeasure}) to \(L_z\) and \(m = n_z - 1\).  The existence of a measure for a ramified critical pair is \tl{Data.CriticalPair.exists\_isTwistMeasure}, which also specifies \(\ms_P = \alg_P\) for \(z \ge 2\) at rank \(2z\).  It uses \tl{Data.CriticalPair.isTwistMeasure\_integralAlgebraicPart} for every critical pair with \(z \ge 2\) and \tl{Data.CriticalPair.exists\_isTwistMeasure\_of\_z\_eq\_one} for ramified pairs with \(z = 1\); the fixed-modulus calculation of Problem~\ref{pr:measure-to-st} is used in \tl{Selmer.exists\_groupRingEval\_eq\_of\_isIntegralMeasure}.  That lemma gives one family \(\St_Q\) with \(1(\St_\varnothing) = \iota(\ms_P(1))\).  At an admissible \(Q\), it interpolates every tame character whose tame part is trivial outside \(Q\) and whose inverse avatar is a character for \(Q\), evaluating the signed character with modulus \(Q\).
\end{leanremark}

\section*{Hints}

\begin{description}
\item[\ref{pr:clq}] Compute \(u/\bar u\) on \(\Z[\ii]^\times\) and put \(\mu = \bar\lambda\); for (b), \(\Z[\ii]\) has class number \(1\) and \(\ii = (1+\ii)/(1-\ii)\).
\item[\ref{pr:insert}] Sieve out multiples of \(\pi_v\) and \(\bar\pi_v\); take \(\pi = 2 + \ii, 3 + 2\ii\).
\item[\ref{pr:weierstrass}] Differentiate the partial-fraction series twice and \(\wp'^2 = 4\wp^3 - 4\wp\) once.  Apply \cite[Thm.~VII.3.4]{Si1}, noting \(2\lfloor e/(q-1)\rfloor \le e\) for \(q \ge 5\).
\item[\ref{pr:poly}] Use Problem~\ref{ec-pr:square} to describe the Galois action on the free rank-one \(\Z[\ii]/q\)-module of \(q\)-torsion.  Conjugation inverts the residue label, and multiplying the torsion parameter by a unit modulo \(q\) translates it.  Apply Problem~\ref{pr:weierstrass}\textup{(b)} to the elementary symmetric coefficients.

At \(q=5\), apply the duplication formula twice in \(x(4P)=x(P)\):
\[
  x(2P) = \frac{(x(P)^2+1)^2}{4x(P)(x(P)^2-1)} .
\]  After discarding \(x(2P)=\pm x(P)\), the equation in \(T=x(P)^2\) is \(T^4-12T^3-26T^2+52T+1=0\).  Its two factors are \(T^2-(6\pm4\sqrt5)T+9\pm4\sqrt5\), with matching signs.  Separate their roots by sign.  The two residue classes have representatives \(1,2\) and \(1+\ii,2(1+\ii)\), respectively, up to Gaussian units.  Count the four residues in each unit orbit when applying Problem~\ref{pr:weierstrass}\textup{(a)}.

For \(q=13,17\), use interval arithmetic on the defining lattice sums.  For each unit residue \(a\) modulo \(q\), choose \(z=a/q\) with \(\lvert\operatorname{Im} z\rvert\le 1/2\).  Sum the rows \(F(z+n\ii)\) for \(-24\le n\le24\), where \(F(z)=\pi^4(2+\cos(2\pi z))/(3\sin^4(\pi z))\) is \(\sum_{m\in\Z}(z+m)^{-4}\).  Use \(\lvert F(z)\rvert\le81\pi^4 e^{-2\pi\lvert\operatorname{Im}z\rvert}\) for \(\lvert\operatorname{Im}z\rvert\ge1\) to bound the omitted rows by a geometric series.  Use the same rows at \(z=0\), replacing the singular row by \(\pi^4/45\), for \(\Omega^4/15\).  Form the polynomial with interval coefficients and isolate the unique integer in each interval scaled by \(q^{6k}\).
The following SageMath code \cite{Sage} encloses the lattice sums needed for the computation.  It can be run with \tl{sage -python}.
\begingroup
\small
\begin{verbatim}
from sage.all import ComplexIntervalField

C = ComplexIntervalField(240)
R, i = C.real_field(), C.gen()
p = R.pi()
tail = 162*p**4*(-49*p).exp() / (1-(-2*p).exp())
e = tail.upper()
error = C(R(-e, e), R(-e, e))

def row(z):
    return p**4*(2+(2*p*z).cos()) / (3*(p*z).sin()**4)

def column(z, omit_zero=False):
    return error + sum(p**4/45 if omit_zero and n == 0
                       else row(z+n*i) for n in range(-24, 25))

G4 = column(C(0), True).real()
\end{verbatim}
\endgroup
\item[\ref{pr:unram-vals}] Write the action of a Galois element on the coefficient labels as \(u \mapsto v u^s\), with \(v \in \Cl_q\) and \(s = \pm 1\), and its action on \(d\)-th roots of unity as raising to the power \(b\), with \((b, d) = 1\).  Then \(\chi(\St_q)\) goes to \(\chi^{-sb}(v)\,\chi^{sb}(\St_q)\), and the phases cancel in pairs \(\{\chi, \chi^{-1}\}\).  At \(q = 13\), use the interval coefficients from the hint to Problem~\ref{pr:poly} and the two cubic Fourier sums.  Use scale \(13^{12}\) for their product.  The \(13\)-torsion field at \(3\) is unramified by \cite[Prop.~VII.4.1(a)]{Si1}.  Over that field, the polynomial \(X^2+3X+3\) of \(\zeta_3-1\) is still Eisenstein.  Use its two roots to compare the valuations of the two Fourier values.
\item[\ref{pr:price}] Consider complex conjugation on \(\Z[\ii]/(3) \cong \F_9\).
\item[\ref{pr:ram-empty}] Compare \(\St^\nu_\varnothing\) with \(\St_5\) via Problem~\ref{pr:poly}; for (b), proceed as in Problem~\ref{pr:insert}.
\item[\ref{pr:ram-ratios}] Sum over nonzero ideals prime to \(5q\) and pair conjugate primes.  For (b), use Problem~\ref{pr:clq}\textup{(b)} to choose representatives that are units at every finite place.  Check both factors of \(\nu\) on the principal unit \(\ii\).  Divide a uniformizer idele's norm-one part by \(\lambda/\bar\lambda\) before evaluating the characters.  At \(q=13\), repeat Problem~\ref{pr:weierstrass}\textup{(a)} at modulus \(65\) with the quadratic weight at \(w_0\).  Galois acts by \(u_m\mapsto v_m u_m^s\) on each residue label \(u_m\in\Cl_m\) for \(m=5,13\), with \(s=1\) or \(-1\) according as it fixes or conjugates \(\ii\).  The quadratic weight contributes the sign \((v_5\mid w_0)\) to each Fourier value.  Check that its square and the cubic phases for the two inverse characters cancel, giving a rational product.  The torsion points have exact order \(65\), so their coordinates are integral by \cite[Thm.~VII.3.4(a)]{Si1}.  Thus \(65^4\) clears each Fourier value into the algebraic integers and \(65^8\) clears their product into \(\Z\).  At modulus \(65\), use the row sums from the hint to Problem~\ref{pr:poly} and the quadratic weight above.  Recover the product of the two Fourier values with scale \(65^8\), then divide by \(1(\St^\nu_\varnothing)^2\).  Compare the two valuations as in the hint to Problem~\ref{pr:unram-vals}, using \(65\)-torsion at \(3\).
\item[\ref{pr:root-val}] Factor \(X^n - 1\) over \(\Ql\), and check \(1 - \zeta \in \Zl\).
\item[\ref{pr:measure-to-st}] Apply \cite{La} place by place over the maximal real subfield of \(E\), rescaling at unramified nonsplit places so almost all coordinates are units; tame characters have weight \(0\).
\item[\ref{pr:norm-measure}] Recall that \(\langle\cdot, \cdot\rangle\) is conjugate-linear in the first variable.
\item[\ref{pr:z-one}] Write \(\bar a = a^{-1}\), and use Problem~\ref{w2-pr:euler-line} for the valuations.
\end{description}

{\small
\begin{thebibliography}{Sage}
\bibitem[GL]{GL} L.~Sernau, \emph{GRH for Dirichlet \(L\)-functions, in Lean~4}, a shorter rewriting of the Lean proof, text by Claude, \url{https://github.com/lsernau/grh-lean/blob/72fead1/grh/README.md}, commit \texttt{72fead1}, readable edition \texttt{grh/}.
\bibitem[Hi]{Hi} H.~Hida, \emph{\(p\)-Adic Automorphic Forms on Shimura Varieties}, Springer Monogr.\ Math., Springer, 2004.
\bibitem[Ka]{Ka} N.~M.~Katz, \emph{\(p\)-adic \(L\)-functions for CM fields}, Invent.\ Math.\ \textbf{49} (1978), 199--297.
\bibitem[La]{La} S.~Lang, \emph{Algebra}, 3rd ed., Grad.\ Texts in Math.\ 211, Springer, 2002.
\bibitem[Sage]{Sage} The Sage Development Team, \emph{SageMath}, \url{https://www.sagemath.org/}.
\bibitem[Si1]{Si1} J.~H.~Silverman, \emph{The Arithmetic of Elliptic Curves}, 2nd ed., Grad.\ Texts in Math.\ 106, Springer, 2009.
\end{thebibliography}
}

\end{document}
